\documentclass{ximera}
%% Macros shared by every activity. Each activity loads this file with
%% %% Macros shared by every activity. Each activity loads this file with
%% %% Macros shared by every activity. Each activity loads this file with
%% \input{../preamble.tex}, and the xourse file with \input{preamble.tex}.

%% MathML intent for screen readers: \xmunits, \xmabs, \xmpoint and the
%% other commands of TailorSwiftBot.sty. The file is latex/xmintent.sty.
\usepackage{xmintent}

%% The interactive pieces of Ximera: \answer, multipleChoice and \choice,
%% selectAll, hint and the rest. Keep this line: without it those commands
%% are undefined and come out as plain text. The file is
%% latex/ximera-components.sty.
\usepackage{ximera-components}

%% Tables with the tabularray package. The settings rowhead, xmcolhead,
%% caption and summary make their header cells known to screen readers, on
%% the web page and in the tagged PDF.
\usepackage{tabularray}

\newcommand{\RR}{\mathbb{R}}
\newcommand{\dfn}{\textit}
, and the xourse file with \input{preamble.tex}.

%% MathML intent for screen readers: \xmunits, \xmabs, \xmpoint and the
%% other commands of TailorSwiftBot.sty. The file is latex/xmintent.sty.
\usepackage{xmintent}

%% The interactive pieces of Ximera: \answer, multipleChoice and \choice,
%% selectAll, hint and the rest. Keep this line: without it those commands
%% are undefined and come out as plain text. The file is
%% latex/ximera-components.sty.
\usepackage{ximera-components}

%% Tables with the tabularray package. The settings rowhead, xmcolhead,
%% caption and summary make their header cells known to screen readers, on
%% the web page and in the tagged PDF.
\usepackage{tabularray}

\newcommand{\RR}{\mathbb{R}}
\newcommand{\dfn}{\textit}
, and the xourse file with \input{preamble.tex}.

%% MathML intent for screen readers: \xmunits, \xmabs, \xmpoint and the
%% other commands of TailorSwiftBot.sty. The file is latex/xmintent.sty.
\usepackage{xmintent}

%% The interactive pieces of Ximera: \answer, multipleChoice and \choice,
%% selectAll, hint and the rest. Keep this line: without it those commands
%% are undefined and come out as plain text. The file is
%% latex/ximera-components.sty.
\usepackage{ximera-components}

%% Tables with the tabularray package. The settings rowhead, xmcolhead,
%% caption and summary make their header cells known to screen readers, on
%% the web page and in the tagged PDF.
\usepackage{tabularray}

\newcommand{\RR}{\mathbb{R}}
\newcommand{\dfn}{\textit}


\title{Tables}

\begin{document}
\begin{abstract}
Tables whose header cells a screen reader can announce.
\end{abstract}
\maketitle

These tables use the \texttt{tabularray} package. Four settings tell a
screen reader how a table is laid out: \texttt{rowhead} (how many rows at the
top hold column headers), \texttt{xmcolhead} (how many columns at the left
hold row headers), \texttt{caption} and \texttt{summary}.

\section{Column headers and row headers}

\begin{center}
\begin{tblr}{colspec={lcc}, hlines, vlines, rowhead=1, xmcolhead=1,
    caption={Quiz scores}, summary={One row for each student, one column for each quiz.}}
Student & Quiz 1 & Quiz 2 \\
Ada     & 9      & 10     \\
Emmy    & 8      & 7      \\
\end{tblr}
\end{center}

\section{Headers that span columns}

\begin{center}
\begin{tblr}{colspec={lcccc}, hlines, vlines, rowhead=2, xmcolhead=1,
    caption={Values of two functions}}
\SetCell[r=2]{c} $x$ & \SetCell[c=2]{c} $f$ & & \SetCell[c=2]{c} $g$ & \\
      & value & slope & value & slope \\
$0$   & $1$   & $0$   & $0$   & $1$   \\
$1$   & $2$   & $2$   & $1$   & $1$   \\
\end{tblr}
\end{center}

\section{A table with an answer}

\begin{problem}
Fill in the missing value.

\begin{center}
\begin{tblr}{colspec={cc}, hlines, vlines, rowhead=1, caption={Squares}}
$n$ & $n^2$ \\
$2$ & $4$ \\
$3$ & $\answer{9}$ \\
\end{tblr}
\end{center}
\end{problem}

\end{document}
